% Otsu 教学例：40 个像素、8 级灰度，计数 2,6,9,5,2,4,8,4。上：直方图与选出的阈值
% k*=3；下：各候选阈值的类间方差 sigma_B^2(k)（由正文公式精确计算后保留两位）。宽约 84mm。
\begin{tikzpicture}[x=9mm,y=2.4mm,font=\figurefont\small]
  \foreach \y in {3,6,9}{\draw[FigureGrid,line width=.4pt] (-.6,\y)--(7.6,\y);
    \node[anchor=east,font=\figurefont\footnotesize] at (-.65,\y) {\y};}
  \foreach \v [count=\k from 0] in {2,6,9,5,2,4,8,4}{
    \ifnum\k<4
      \fill[FigureBlueFill,draw=FigureBlue,line width=.6pt] (\k-.35,0) rectangle (\k+.35,\v);
    \else
      \fill[FigureYellowFill,draw=FigureYellow,line width=.6pt] (\k-.35,0) rectangle (\k+.35,\v);
    \fi
    \node[below,font=\figurefont\footnotesize] at (\k,0) {\k};
  }
  \draw[FigureStroke,line width=.6pt] (-.6,0)--(7.6,0);
  \draw[FigureRed,dashed,line width=.8pt] (3.5,-.5)--(3.5,11);
  \node[anchor=south,font=\figurefont\footnotesize] at (3.5,11) {$k^*=3$};
  \node[font=\figurefont\footnotesize,text=FigureMuted] at (1.5,11.5) {$C_0$};
  \node[font=\figurefont\footnotesize,text=FigureMuted] at (5.5,11.5) {$C_1$};
  \node[anchor=east] at (-1.4,5) {像素数};
  \node at (3.5,-3.4) {灰度级 $i$};
  % sigma_B^2 curve
  \begin{scope}[yshift=-12mm,y=2.4mm]
    \draw[FigureGrid,line width=.4pt] (-.6,0)--(7.6,0);
    \foreach \y in {2,4}{\draw[FigureGrid,line width=.4pt] (-.6,{-6+\y})--(7.6,{-6+\y});
      \node[anchor=east,font=\figurefont\footnotesize] at (-.65,{-6+\y}) {\y};}
    \draw[FigureStroke,line width=.6pt] (-.6,-6)--(7.6,-6);
    \draw[FigureBlue,line width=1pt] (0,{-6+0.67})--(1,{-6+2.00})--(2,{-6+3.46})--(3,{-6+3.97})--(4,{-6+3.92})--(5,{-6+3.26})--(6,{-6+1.30});
    \foreach \k/\v in {0/0.67,1/2.00,2/3.46,3/3.97,4/3.92,5/3.26,6/1.30}{
      \fill[FigureBlue] (\k,{-6+\v}) circle (.7mm);}
    \node[anchor=west,font=\figurefont\footnotesize] at (3.1,{-6+4.5}) {$3.97$};
    \foreach \k in {0,...,6}{\node[below,font=\figurefont\footnotesize] at (\k,-6) {\k};}
    \node[anchor=east] at (-1.4,-3) {$\sigma_B^2(k)$};
    \node at (3.5,-9.4) {候选阈值 $k$};
  \end{scope}
\end{tikzpicture}
